\setcounter{reportfigure}{0}
\renewcommand{\thereportfigure}{A.\arabic{reportfigure}}
\section*{Appendix A --- Part I: Simulink diagrams}
\subsection*{Question 1: inverted pendulum}
\begin{center}
\includegraphics[width=\linewidth,height=0.25\textheight,keepaspectratio]{pendulum_plant_interior.png}
\reportcaption{Opened pendulum plant: angular acceleration and the two state integrators.}
\end{center}
\begin{center}
\includegraphics[width=\linewidth,height=0.25\textheight,keepaspectratio]{pendulum_regulation_diagram.png}
\reportcaption{Simulink diagram export of the pendulum closed loop, including state feedback, nonlinear cancellation, plant and observation scopes.}
\end{center}
\begin{center}
\includegraphics[width=0.92\linewidth,height=0.25\textheight,keepaspectratio]{pendulum_linearization_controller.png}
\reportcaption{Simulink diagram export of the pendulum's nonlinear cancellation law.}
\end{center}
\clearpage
\subsection*{Question 2: DC motor}
\begin{center}
\includegraphics[width=\linewidth,height=0.25\textheight,keepaspectratio]{dc_motor_plant_interior.png}
\reportcaption{Opened DC motor plant: electrical and mechanical dynamics with three state integrators.}
\end{center}
\begin{center}
\includegraphics[width=\linewidth,height=0.25\textheight,keepaspectratio]{dc_motor_tracking_diagram.png}
\reportcaption{Simulink diagram export of the motor position-tracking loop. Position, speed and current feed the linear tracking control and nonlinear voltage calculation.}
\end{center}
\begin{center}
\includegraphics[width=\linewidth,height=0.25\textheight,keepaspectratio]{dc_motor_linearizing_controller.png}
\reportcaption{Simulink diagram export of the motor cancellation controller. Acceleration is calculated from speed and current.}
\end{center}
\clearpage
\subsection*{Question 3: non-isothermal CSTR}
\begin{center}
\includegraphics[width=\linewidth,height=0.25\textheight,keepaspectratio]{cstr_plant_interior.png}
\reportcaption{Opened CSTR plant: mass and energy balances with concentration and temperature integrators.}
\end{center}
\begin{center}
\includegraphics[width=\linewidth,height=0.25\textheight,keepaspectratio]{cstr_tracking_diagram.png}
\reportcaption{Simulink diagram export of the reactor temperature-tracking loop. The error gain commands a virtual temperature rate.}
\end{center}
\begin{center}
\includegraphics[width=\linewidth,height=0.25\textheight,keepaspectratio]{cstr_linearizing_controller.png}
\reportcaption{Simulink diagram export of the coolant-temperature calculation, including conversion from minute-based process rates to seconds.}
\end{center}
\clearpage
\subsection*{Question 4: ball and beam}
\begin{center}
\includegraphics[width=\linewidth,height=0.25\textheight,keepaspectratio]{ballbeam_plant_interior.png}
\reportcaption{Opened ball-and-beam plant: ball and beam accelerations and the four state integrators.}
\end{center}
\begin{center}
\includegraphics[width=\linewidth,height=0.45\textheight,keepaspectratio]{ballbeam_closed_loop.png}
\reportcaption{Diagnostic Simulink model of the ball-and-beam closed loop. The linear controller acts on the new coordinates, and the designed linear closed loop runs alongside for comparison.}
\end{center}
\clearpage
\subsection*{Question 4: linearizing law and coordinates}
\begin{center}
\includegraphics[width=\linewidth,height=0.25\textheight,keepaspectratio]{ballbeam_linearizing_law.png}
\reportcaption{Simulink diagram export of the ball-and-beam linearizing law $u=\alpha(x)+\beta(x)v$.}
\end{center}
\begin{center}
\includegraphics[width=\linewidth,height=0.25\textheight,keepaspectratio]{ballbeam_coordinates.png}
\reportcaption{Simulink diagram export of the new coordinates $z$: ball position and velocity, then acceleration and jerk with the centrifugal term neglected.}
\end{center}
